% GENERATED FILE. Edit canonical lessons, then run npm run generate:books.
% canonical-source-hash: fnv1a64:71431b0396700da4
% canonical-lessons: PT-C299-acidente, PT-C299-choque, PT-C299-chocar, PT-C299-colidir, PT-C299-sinistro

\chapter{Um acidente na estrada}
\label{ch:pt-um-acidente-na-estrada}
\hlchaptermodality{299}

\begin{chapteropening}
\textbf{By the end of this chapter:} \emph{I can describe a road accident: the crash, what collided with what, and the mishap.}

Complete the last lesson of chapter 299: I can describe a road accident: the crash, what collided with what, and the mishap.
\end{chapteropening}

\section[o acidente]{o acidente — the accident}

\label{lesson:PT-C299-acidente}



\begin{quote}
Before the new one: say the Portuguese for the inmate, then the Portuguese for to set free.
\end{quote}

\subsection*{You'll want to know: o acidente}
\textbf{o acidente} — \textquotedblleft{}the accident\textquotedblright{}.

\begin{quote}
Houve um acidente na autoestrada.
\end{quote}

\begin{quote}
\emph{There was an accident on the motorway.}
\end{quote}

\emph{Ter um acidente} — to have an accident; \emph{por acidente} — by chance.

\begin{cousinweb}
\emph{acidente} is a learned borrowing from Latin \emph{accidēns} \textquotedblleft{}happening, chance\textquotedblright{} (EtymDB), from \emph{accidere} \textquotedblleft{}to fall upon, to happen\textquotedblright{}, built on \emph{cadere} \textquotedblleft{}to fall\textquotedblright{}.
\end{cousinweb}

\begin{grammarlens}[title={What you've built}]
One more word for telling, explaining and describing at B1.
\end{grammarlens}

\subsection*{Your turn}
\begin{itemize}
\raggedright
  \item \emph{Say it:} \emph{o acidente}
  \item \emph{Say it:} \emph{o acidente}, once more
  \item \emph{Say it:} a sentence of your own with \emph{o acidente}
  \item \emph{Recall:} say the Portuguese for the inmate, then the Portuguese for to set free, then say \emph{o acidente} again
\end{itemize}

\begin{tcolorbox}[breakable,colback=teal!4,colframe=teal!35!black,title={Before you move on}]
What does o acidente mean? (\textquotedblleft{}The accident\textquotedblright{}.) Say it once more.
\end{tcolorbox}
\section[o choque]{o choque — the crash, the impact; the shock}

\label{lesson:PT-C299-choque}



\begin{quote}
Before the new one: say the Portuguese for to set free, then the Portuguese for the accident.
\end{quote}

\subsection*{You'll want to know: o choque}
\textbf{o choque} — \textquotedblleft{}the crash, the impact; the shock\textquotedblright{}.

\begin{quote}
O choque entre os dois carros foi forte.
\end{quote}

\begin{quote}
\emph{The crash between the two cars was violent.}
\end{quote}

Also an electric shock or an emotional one: \emph{ficar em choque}.

\begin{cousinweb}
\emph{choque} was borrowed from French \emph{choc} \textquotedblleft{}a blow, an impact\textquotedblright{} (EtymDB).
\end{cousinweb}

\begin{grammarlens}[title={What you've built}]
One more word for telling, explaining and describing at B1.
\end{grammarlens}

\subsection*{Your turn}
\begin{itemize}
\raggedright
  \item \emph{Say it:} \emph{o choque}
  \item \emph{Say it:} \emph{o choque}, once more
  \item \emph{Say it:} a sentence of your own with \emph{o choque}
  \item \emph{Recall:} say the Portuguese for to set free, then the Portuguese for the accident, then say \emph{o choque} again
\end{itemize}

\begin{tcolorbox}[breakable,colback=teal!4,colframe=teal!35!black,title={Before you move on}]
What does o choque mean? (\textquotedblleft{}The crash, the impact; the shock\textquotedblright{}.) Say it once more.
\end{tcolorbox}
\section[chocar]{chocar — to crash, to run into; to shock}

\label{lesson:PT-C299-chocar}



\begin{quote}
Before the new one: say the Portuguese for the accident, then the Portuguese for the crash.
\end{quote}

\subsection*{You'll want to know: chocar}
\textbf{chocar} — \textquotedblleft{}to crash, to run into; to shock\textquotedblright{}.

\begin{quote}
O autocarro chocou contra um muro.
\end{quote}

\begin{quote}
\emph{The bus crashed into a wall.}
\end{quote}

\emph{Chocar com} or \emph{chocar contra}. Also ‘to shock’: \emph{a notícia chocou o país}.

\begin{cousinweb}
\emph{chocar} is built on \emph{choque} \textquotedblleft{}an impact\textquotedblright{} (EtymDB), from French \emph{choc}, with the verb ending \emph{-ar}.
\end{cousinweb}

\begin{grammarlens}[title={What you've built}]
One more word for telling, explaining and describing at B1.
\end{grammarlens}

\subsection*{Your turn}
\begin{itemize}
\raggedright
  \item \emph{Say it:} \emph{chocar}
  \item \emph{Say it:} \emph{chocar}, once more
  \item \emph{Say it:} a sentence of your own with \emph{chocar}
  \item \emph{Recall:} say the Portuguese for the accident, then the Portuguese for the crash, then say \emph{chocar} again
\end{itemize}

\begin{tcolorbox}[breakable,colback=teal!4,colframe=teal!35!black,title={Before you move on}]
What does chocar mean? (\textquotedblleft{}To crash, to run into; to shock\textquotedblright{}.) Say it once more.
\end{tcolorbox}
\section[colidir]{colidir — to collide}

\label{lesson:PT-C299-colidir}



\begin{quote}
Before the new one: say the Portuguese for the crash, then the Portuguese for to crash.
\end{quote}

\subsection*{You'll want to know: colidir}
\textbf{colidir} — \textquotedblleft{}to collide\textquotedblright{}.

\begin{quote}
Dois camiões colidiram na ponte.
\end{quote}

\begin{quote}
\emph{Two lorries collided on the bridge.}
\end{quote}

Formal, a news word; in speech, \emph{bater} or \emph{chocar}.

\begin{cousinweb}
\emph{colidir} is a learned borrowing from Latin \emph{collīdere} \textquotedblleft{}to strike together\textquotedblright{}, built on \emph{laedere} \textquotedblleft{}to strike, to hurt\textquotedblright{} (Lewis \& Short).
\end{cousinweb}

\begin{grammarlens}[title={What you've built}]
One more word for telling, explaining and describing at B1.
\end{grammarlens}

\subsection*{Your turn}
\begin{itemize}
\raggedright
  \item \emph{Say it:} \emph{colidir}
  \item \emph{Say it:} \emph{colidir}, once more
  \item \emph{Say it:} a sentence of your own with \emph{colidir}
  \item \emph{Recall:} say the Portuguese for the crash, then the Portuguese for to crash, then say \emph{colidir} again
\end{itemize}

\begin{tcolorbox}[breakable,colback=teal!4,colframe=teal!35!black,title={Before you move on}]
What does colidir mean? (\textquotedblleft{}To collide\textquotedblright{}.) Say it once more.
\end{tcolorbox}
\section[o sinistro]{o sinistro — the mishap; an accident (in insurance language)}

\label{lesson:PT-C299-sinistro}



\begin{quote}
Before the new one: say the Portuguese for to crash, then the Portuguese for to collide.
\end{quote}

\subsection*{You'll want to know: o sinistro}
\textbf{o sinistro} — \textquotedblleft{}the mishap; an accident (in insurance language)\textquotedblright{}.

\begin{quote}
O seguro pagou o sinistro.
\end{quote}

\begin{quote}
\emph{The insurance paid for the accident.}
\end{quote}

An insurance and news word. As an adjective, \emph{sinistro} means ‘sinister’.

\begin{cousinweb}
\emph{sinistro} was borrowed from Latin \emph{sinister} \textquotedblleft{}on the left; unlucky\textquotedblright{} (EtymDB): the left side was the side of bad omens.
\end{cousinweb}

\begin{grammarlens}[title={What you've built}]
One more word for telling, explaining and describing at B1.
\end{grammarlens}

\subsection*{Your turn}
\begin{itemize}
\raggedright
  \item \emph{Say it:} \emph{o sinistro}
  \item \emph{Say it:} \emph{o sinistro}, once more
  \item \emph{Say it:} a sentence of your own with \emph{o sinistro}
  \item \emph{Recall:} say the Portuguese for to crash, then the Portuguese for to collide, then say \emph{o sinistro} again
\end{itemize}

\begin{tcolorbox}[breakable,colback=teal!4,colframe=teal!35!black,title={Before you move on}]
What does o sinistro mean? (\textquotedblleft{}The mishap; an accident (in insurance language)\textquotedblright{}.) Say it once more.
\end{tcolorbox}
